%% ============================================================
%% sections/a5-evidence.tex
%% H. R. 9510 Bill v3.0 - APPENDIX A. VISUAL ENGINEERING EVIDENCE (NON-OPERATIVE).
%% The visualization of Bill v2.0: the VVUQ-01 and VVUQ-02 engineering record that
%% grounds the findings. Rendered visuals: Figure 6 (gate funnel comparison),
%% Figure 7 (procedure timelines), Figure 8 (humanoid safety figure), Table 5
%% (external standards map).
%% ============================================================
% \clearpage
\phantomsection\label{toc:a5}
\addcontentsline{toc}{section}{Appendix A. Visual Engineering Evidence (Non-Operative)}
\usccrosshead{APPENDIX A. VISUAL ENGINEERING EVIDENCE (NON-OPERATIVE)}

\uscprov{1}{This appendix is non-operative. It visualizes the engineering record
of VVUQ-01 and VVUQ-02 that grounds the findings in section 2 and the gate
discipline codified in section 515D, so a reader can see the evidence the bill
rests on without leaving the document. None of these figures is cited as the
legal basis of an operative clause; only the enacted statutes, in-force rules, and
recognized standards named in the bill carry operative weight
\cite{kawchak2026vvuq01, kawchak2026vvuq02}.}

\uscprov{1}{The gate is the load-bearing object of the whole record. VVUQ-01 holds
it to a single behavior across three dimensions, accepting one candidate and
blocking five; VVUQ-02 scales it to ten humanoid-specific gates without changing
the ACCEPT, BLOCK, ESCALATE semantics. Figure 6 sets the two side by side, and the
verification fraction stays an equality at 1.0 in both, which is why the bill can
require it as an equality \cite{kawchak2026vvuq01, kawchak2026vvuq02}.}

\begin{asciifig}
VVUQ-01: 6 candidates, 3 dimensions   VVUQ-02: 10 humanoid gates
+---------------------------------+   +---------------------------------+
| verify  : fraction == 1.0       |   | [01] handeye-servo    V>=0.97   |
| validate: agree>=0.95 err<=0.05 |   | [02] finger-force     V>=0.95   |
| quantify: CV<=0.10 over >=3 runs|   | [03] whole-body-bal   V>=0.98   |
+---------------------------------+   | [04] plan-correct     V>=0.95   |
 case 1 clean ............ ACCEPT     | [05] grasp-handover   V>=0.96   |
 case 2 no human review .. BLOCK      | [06] vascular-no-fly* V=1.00    |
 case 3 high CV 0.365 .... ESCALATE   | [07] suturing         V>=0.96   |
 case 4 disagree ref ..... BLOCK      | [08] perception       V>=0.95   |
 case 5 only 2 runs ...... BLOCK      | [09] human-collision* V=1.00    |
 case 6 over 200K cap .... BLOCK      | [10] fault-estop    * V=1.00    |
                 |                    +---------------------------------+
                 v                                      v
 1 ACCEPT, 5 BLOCK (1 escalates)       ACCEPT iff all 10; composite 93.56
 verification fraction held at 1.0     * three hard-catastrophe gates
\end{asciifig}
\vspace{-0.7cm}
\figcaption{Figure 6. The gate, from one dimension to ten. The VVUQ-01
six-candidate funnel that accepted one\\artifact and blocked five, beside the
VVUQ-02 ten-gate funnel that cleared a surgical humanoid\\behavior against ten
standards-bound gates, three of them carrying a hard catastrophe predicate.}

\uscprov{1}{Figure 7 reads the assurance windows on one page: the
sixty-second eight-phase robotic Whipple a single humanoid performs with two
hands, over the one-hundred-sixty-eight-day pancreatic-cancer adjuvant pilot the
procedure opens \cite{kawchak2026vvuq02, kawchak2026vvuq01}.}

\begin{asciifig}
60-Second 8-Phase Whipple Timeline (single humanoid, two hands)

  s  0      6        16      24      32          42      48     54     60
     |------|--------|-------|-------|-----------|-------|------|------|
  P1 |Kocher|
  P2        |venous  |
            |control |
  P3                 |uncin. |
  P4                         |specimen|
  P5                                  |PJ suture  |
  P6                                              |HJ     |
  P7                                                      |GJ    |
  P8                                                             |close|

168-Day PDAC Hybrid Pilot (the same procedure, then adjuvant therapy)

  day 0        28        56        84       112       140       168
   |---------|---------|---------|---------|---------|---------|
   [Whipple]  [cyc 1]   [cyc 2]   [cyc 3]   [cyc 4]   [cyc 5]   [cyc 6]
   8-arm      adjuvant Daraxonrasib, one cycle every 28 days
   60 s
\end{asciifig}
\vspace{-0.7cm}
\figcaption{Figure 7. The procedure timelines: the sixty-second eight-phase
Whipple window over the one-hundred-\\sixty-eight-day adjuvant pilot, so the
single-procedure assurance window and the full pilot window read together.}

\uscprov{1}{The safety figure is where embodiment becomes visible. Figure 8
combines the seventy-one-degree-of-freedom kinematic chain, the balance
zero-moment-point and support polygon, and the vessel no-fly proximity bands: the
kinematic chain shows how many ways a body can move, the support polygon shows the
balance margin a fall would cross, and the no-fly bands show the vascular volume
the hard predicate forbids breaching. These are the dimensions a wrong motion
would violate and a human could not catch in time \cite{isots15066, iec8060127}.}

\begin{asciifigsmall}
(A) H2-Surgical Kinematic Chain (71 DOF, hypothetical 2030)

                          [head pan/tilt] 2 DOF
                                   |
                               [ head ]
                                   |
              [L arm 7 DOF]----[ torso ]----[R arm 7 DOF]
                    |          [waist 3]          |
              [L hand 20 DOF]      |        [R hand 20 DOF]
              5 fingers x 4        |        5 fingers x 4
                              [ pelvis ]
                              /        \
                      [L leg 6]        [R leg 6]
                           |                |
                      [L foot]          [R foot]
                  ==========  support polygon  ==========
              arms 14 + hands 40 + stance 15 + head 2 = 71 DOF

(B) Balance ZMP and Support Polygon (VVUQ 03, top-down)

   +175 +-------------------------------+  <- support polygon edge
        |        warning ring (15 mm)   |
        |      .......................  |
        |      :  floor ring (8 mm)  :  |
      0 |      :      +-----+        :  |  ZMP margin 45 mm (nominal)
        |      :      | ZMP |  o CoM :  |  ACCEPT while margin >= 8 mm
        |      :.....................:  |
   -175 +-------------------------------+
        -130             0           +130  x (mm)

(C) Vessel No-Fly Proximity Bands (VVUQ 06, cross-section)

   tip far .................................. clear     v=1.0  cap 3.0 N
   | 6 mm ----------------------------------- no_fly    v=0.5  cap 2.5 N
   | 4 mm ----------------------------------- soft_warn v=0.1  cap 1.5 N
   | 2 mm =================================== HARD STOP v=0.0  e-stop
   o vessel centerline (SMV, PV, hepatic, celiac, SMA); 0 breaches required
\end{asciifigsmall}
\vspace{-0.7cm}
\figcaption{Figure 8. The humanoid safety figure: the kinematic chain, the
balance zero-moment-point and support polygon, and the vessel no-fly proximity
bands, the three drawn artifacts that make the catastrophe gates concrete.}

\uscprov{1}{Table 5 maps the recognized consensus standards bound to the gates and
duties, so each row of the threshold schedule in section 515D traces to a
published standard \cite{asmevv40, isots15066, iec60601, iso13849}.}

{\footnotesize
\begin{xltabular}{\textwidth}{@{}L{3cm} Y L{4.5cm}@{}}
\hline\noalign{\vskip 0.04cm}
\textbf{Standard} & \textbf{What it governs} & \textbf{Gate or duty}\\
\hline\noalign{\vskip 0.04cm}
\endfirsthead
\hline\noalign{\vskip 0.04cm}
\textbf{Standard} & \textbf{What it governs} & \textbf{Gate or duty}\\
\hline\noalign{\vskip 0.04cm}
\endhead
\hline\noalign{\vskip 0.04cm}
\endlastfoot
ASME V\&V 40-2018 & Model credibility (risk equals influence times consequence) &
Gates 01, 08; gate credibility basis\\
\hline\noalign{\vskip 0.04cm}
NASA-STD-7009A & Uncertainty quantification as a credibility factor & The
coefficient-of-variation bounds\\
\hline\noalign{\vskip 0.04cm}
IEEE Std 7009-2024 & Fail-safe design and hand-back-to-human & Gates 04, 10; the
ESCALATE rule\\
\hline\noalign{\vskip 0.04cm}
IEC 80601-2-77 & Basic safety of robotically assisted surgery & Gates 01, 02, 05,
06, 07\\
\hline\noalign{\vskip 0.04cm}
ISO 14971 & Risk management; risk control verification & Gate 06; hard
predicate: risk limit\\
\hline\noalign{\vskip 0.04cm}
ISO/TS 15066 & Power-and-force limits for human contact & Gates 02, 09\\
\hline\noalign{\vskip 0.04cm}
ISO 13482 & Stability and collision for close-contact robots & Gates 03, 09\\
\hline\noalign{\vskip 0.04cm}
IEC 60601-1 & Single-fault safe state for medical electrical equipment & Gates 03,
10; cybersecur., lifecycle\\
\hline\noalign{\vskip 0.04cm}
ISO 13849-1 & Performance level of the safety control path & Gate 10; the
emergency-stop path\\
\hline\noalign{\vskip 0.04cm}
ISO 9283 & Pose accuracy and repeatability & Gates 01, 05\\
\hline\noalign{\vskip 0.04cm}
21 CFR part 11 & Electronic-record integrity & Hash-chained trail, subsection (e)\\
\hline
\end{xltabular}}
\vspace{-0.8cm}
\figcaption{Table 5. External standards map: the recognized consensus standards
bound\\to the gates and duties of section 515D, with the gate or duty each
governs.}
